\section{Model Serangan Berdasarkan Informasi Penyerang}

Kekuatan seorang kriptanalis ditentukan oleh jumlah dan jenis informasi yang ia miliki. Dalam studi kriptografi, kita menggunakan model serangan untuk mensimulasikan skenario terburuk yang mungkin dihadapi oleh sebuah sistem. Semakin banyak informasi yang dapat diakses oleh penyerang, semakin rendah tingkat kesulitan untuk memecahkan cipher tersebut.

\subsection{Ciphertext-Only Attack (COA)}

\textit{Ciphertext-Only Attack} adalah model serangan yang paling menantang karena penyerang \textbf{hanya memiliki akses ke satu atau lebih cipherteks} yang berhasil dicegat dari saluran komunikasi. Penyerang tidak mengetahui isi plainteks maupun kunci yang digunakan.

\textbf{Strategi Utama}:
\begin{itemize}
    \item \textbf{Analisis Statistik}: Menghitung frekuensi huruf, bigram, atau trigram untuk menebak bahasa dan struktur pesan.
    \item \textbf{Brute Force}: Mencoba semua kemungkinan kunci jika ruang kuncinya kecil (misal: Caesar Cipher).
    \item \textbf{Analisis Pola}: Mencari pengulangan dalam cipherteks untuk menentukan panjang kunci (seperti Metode Kasiski pada Vigenère).
\end{itemize}
COA adalah model standar bagi penyadap yang hanya mampu menguping saluran komunikasi tanpa memiliki kontrol atas pengiriman pesan.

\subsection{Known-Plaintext Attack (KPA)}

Dalam \textit{Known-Plaintext Attack}, penyerang memiliki satu atau lebih pasangan \textbf{plainteks dan cipherteks yang bersesuaian}.

\textbf{Sumber Plainteks yang Diketahui}:
\begin{itemize}
    \item \textbf{Standard Headers}: Banyak protokol komunikasi memiliki format tetap. Misalnya, setiap dokumen PDF selalu dimulai dengan karakter \texttt{"\%PDF"}, atau email bisnis sering dimulai dengan \texttt{"Dengan hormat,"}.
    \item \textbf{Tebakan Logis}: Penyerang mungkin tahu bahwa pesan tersebut mengandung kata-kata tertentu berdasarkan konteks (misal: \texttt{"Laporan Tahunan"} atau \texttt{"Password:"}).
    \item \textbf{Kebocoran Parsial}: Sebagian isi pesan bocor melalui saluran lain, memberikan petunjuk tentang hubungan antara plainteks dan cipherteks.
\end{itemize}
KPA sangat mematikan bagi cipher linear seperti Hill Cipher, di mana pasangan $(\mathbf{P}, \mathbf{C})$ dapat digunakan untuk menghitung matriks kunci melalui operasi invers.

\subsection{Chosen-Plaintext Attack (CPA)}

Pada \textit{Chosen-Plaintext Attack}, penyerang memiliki kemampuan untuk \textbf{memilih plainteks apa pun dan mendapatkan cipherteks yang bersesuaian}. Penyerang tidak mengetahui kuncinya, tetapi memiliki akses ke "kotak hitam" (perangkat enkripsi) yang akan mengenkripsi pesan apa pun yang ia masukkan.

\textbf{Strategi Utama}:
Penyerang dapat memasukkan pola data yang sangat spesifik (misalnya, string berisi semua huruf 'A' atau pola bit $010101$) untuk mengamati bagaimana algoritma merespons. Hal ini memungkinkan penyerang untuk memetakan struktur internal cipher atau mempersempit ruang kunci secara drastis.

\subsection{Chosen-Ciphertext Attack (CCA)}

\textit{Chosen-Ciphertext Attack} adalah model serangan yang paling kuat. Penyerang dapat \textbf{memilih cipherteks mana pun dan mendapatkan hasil dekripsinya (plainteks)}.

\textbf{Mekanisme}:
Penyerang mengirimkan cipherteks yang telah dimodifikasi ke sistem dekripsi (seperti server yang memberikan respons error atau konfirmasi) dan menganalisis plainteks yang keluar. Informasi ini kemudian digunakan untuk menyimpulkan kunci rahasia.

\textbf{Adaptive CCA (CCA2)}: Variasi di mana penyerang memilih cipherteks berikutnya \textit{berdasarkan} hasil dekripsi dari cipherteks sebelumnya. Ini memungkinkan penyerang untuk melakukan "interogasi" terhadap sistem dekripsi secara iteratif hingga kunci terungkap.

\begin{table}[htbp]
\centering
\small
\begin{tabular}{|l|p{6cm}|l|}
\hline
\textbf{Model Serangan} & \textbf{Informasi yang Dimiliki Penyerang} & \textbf{Tingkat Kesulitan} \\ \hline
COA & Hanya satu atau lebih cipherteks & Sangat Tinggi \\ \hline
KPA & Pasangan (Plainteks, Cipherteks) & Tinggi \\ \hline
CPA & Pilih Plainteks $\to$ Dapat Cipherteks & Sedang \\ \hline
CCA & Pilih Cipherteks $\to$ Dapat Plainteks & Rendah \\ \hline
\end{tabular}
\caption{Perbandingan model serangan berdasarkan informasi yang tersedia.}
\end{table}
